\chapter{复杂航天器 (Complex Spacecraft)}

相比于简单航天器，复杂航天器的动力学方程推导往往异常繁琐，需要采用特别高效的方法。本章前七节建立了这一方法论，并在后续四节中将其应用于完全或部分由弹性组件构成的航天器。

\section{广义主动力 (Generalized Active Forces)}

\begin{mydefinition}{广义主动力的定义}
给定一个由 \(N\) 个粒子 \(P_1, \dots, P_N\) 组成的系统 \(S\)，假设定义了 \(n\) 个广义速度（参见 1.21 节）。设 \(\mathbf{v}_r^{P_i}\) 为 \(P_i\) 在参考系 \(A\) 中的第 \(r\) 个偏速度，\(\mathbf{R}_i\) 为作用在 \(P_i\) 上的所有接触力和体积力的合力。那么，\(S\) 在 \(A\) 中的广义主动力 \(F_1, \dots, F_n\) 定义为：
\begin{equation}
F_r = \sum_{i=1}^N \mathbf{v}_r^{P_i} \cdot \mathbf{R}_i \qquad (r=1,\dots,n) \tag{1}
\end{equation}
\end{mydefinition}

并非所有对 \(\mathbf{R}_i\) 有贡献的力都会对 \(F_r\) 产生贡献。例如，刚体光滑表面上粒子之间的接触力对 \(F_r\) 的总贡献为零。此外，若刚体 \(B\) 内部的粒子之间的接触力和引力等效于一个作用在点 \(Q\) 的合力 \(\mathbf{R}\) 和一个力偶矩 \(\mathbf{T}\)，则它们对 \(F_r\) 的贡献可以表示为：
\begin{equation}
(F_r)_B = \boldsymbol{\omega}_r \cdot \mathbf{T} + \mathbf{v}_r^Q \cdot \mathbf{R} \qquad (r=1,\dots,n) \tag{2}
\end{equation}
其中 \(\boldsymbol{\omega}_r\) 和 \(\mathbf{v}_r^Q\) 分别是 \(B\) 在 \(A\) 中的第 \(r\) 个偏角速度和点 \(Q\) 的第 \(r\) 个偏速度。

\begin{proof}[内力与约束力贡献为零的证明]
对于作用在光滑表面上的接触力 \(C = C\mathbf{n}\)，由于 \(P\) 的速度在法线方向上的分量为零（既不能穿透刚体又不能离开），其对广义力的贡献为零。
对于刚体内部粒子间的相互作用力，根据牛顿第三定律，内力大小相等、方向相反且共线。利用偏速度关系 \(\mathbf{v}_r^{P_j} = \mathbf{v}_r^{P_i} + \boldsymbol{\omega}_r \times \mathbf{p}_{ij}\) 进行向量点积运算，可证明内力对广义力的总贡献相互抵消。由力系等效性，即可得到刚体合力和合力矩对广义主动力的贡献公式 (2)。
\end{proof}

\section{势能 (Potential Energy)}

\begin{mydefinition}{势能}
如果系统 \(S\) 在参考系 \(A\) 中的构型由 \(n\) 个广义坐标 \(q_1,\dots,q_n\) 描述，且存在一个函数 \(W(q_1,\dots,q_n,t)\)，使得所有广义主动力 \(F_1,\dots,F_n\) 都可以表示为：
\begin{equation}
F_r = -\frac{\partial W}{\partial q_r} \qquad (r=1,\dots,n) \tag{1}
\end{equation}
那么 \(W\) 称为 \(S\) 在 \(A\) 中的势能。
\end{mydefinition}

在某些情况下，2.10-2.19 节中讨论的力函数可以直接提供势能。然而，可能存在明显的力函数，但势能却不存在的情况。例如，仅用中心引力近似合力，而用二阶项近似力矩时，系统就没有势能。此时，尽管不能使用势能，但该近似在航天器动力学中依然非常有用且广泛采用。

\begin{proof}[重力势能推导]
对于两个粒子相互作用的情形，合力 \(\mathbf{F}\) 正是标量势 \(V\) 的梯度。因此 \(W = -V\) 即为势能。通过构建包含刚体的多体系统的广义主动力，可以证明推广到刚体的情形时，系统广义力 \(F_r = \sum \mathbf{v}_r^{P_i} \cdot \mathbf{F}_i\) 同样可通过 \(W = -V\) 求负梯度得到，从而证明了式 (1) 在重力场中的有效性。
\end{proof}

\section{广义惯性力 (Generalized Inertia Forces)}

\begin{mydefinition}{广义惯性力}
给定系统 \(S\) 由 \(N\) 个粒子组成，设有 \(n\) 个广义速度。设 \(\mathbf{v}_r^{P_i}\) 为 \(P_i\) 在参考系 \(A\) 中的第 \(r\) 个偏速度（1.21节），\(\mathbf{R}_i^*\) 为 \(P_i\) 在 \(A\) 中的惯性力，即：
\begin{equation}
\mathbf{R}_i^* = -m_i \mathbf{a}_i \qquad (i=1,\dots,N) \tag{1}
\end{equation}
那么 \(S\) 在 \(A\) 中的广义惯性力 \(F_1^*, \dots, F_n^*\) 定义为：
\begin{equation}
F_r^* = \sum_{i=1}^N \mathbf{v}_r^{P_i} \cdot \mathbf{R}_i^* \qquad (r=1,\dots,n) \tag{2}
\end{equation}
\end{mydefinition}

对于属于 \(S\) 的刚体 \(B\)，所有粒子惯性力对 \(F_r^*\) 的贡献可表示为：
\begin{equation}
(F_r^*)_B = \boldsymbol{\omega}_r \cdot \mathbf{T}^* + \mathbf{v}_r \cdot \mathbf{R}^* \qquad (r=1,\dots,n) \tag{5}
\end{equation}
其中 \(\mathbf{R}^* = -M\mathbf{a}^*\) 是平动惯性力，\(\mathbf{T}^*\) 是**惯性力矩**：
\begin{equation}
\mathbf{T}^* = -\boldsymbol{\alpha} \cdot \mathbf{I} - \boldsymbol{\omega} \times \mathbf{I} \cdot \boldsymbol{\omega} \tag{6}
\end{equation}
这里的 \(\boldsymbol{\alpha}\) 和 \(\boldsymbol{\omega}\) 分别是 \(B\) 在 \(A\) 中的角加速度和角速度，\(\mathbf{I}\) 是 \(B\) 的中心惯性并矢。

若将惯性力矩投影到与主惯性轴重合的任意单位基 \(c_1,c_2,c_3\) 上，可得：
\begin{equation}
\mathbf{T}^* = -[\alpha_1 I_1 - \omega_2\omega_3(I_2 - I_3)]c_1 - [\alpha_2 I_2 - \omega_3\omega_1(I_3 - I_1)]c_2 - [\alpha_3 I_3 - \omega_1\omega_2(I_1 - I_2)]c_3 \tag{8}
\end{equation}

\begin{proof}[惯性力矩 \(T^*\) 的推导]
根据定义 \(\mathbf{T}^* = -\sum m_i \mathbf{r}_i \times \mathbf{a}_i\)，将加速度 \(\mathbf{a}_i = \mathbf{a}^* + \boldsymbol{\alpha} \times \mathbf{r}_i + \boldsymbol{\omega} \times (\boldsymbol{\omega} \times \mathbf{r}_i)\) 代入，利用质心定义 \(\sum m_i \mathbf{r}_i = 0\) 和惯性并矢 \(\mathbf{I} = \sum m_i(\mathbf{r}_i^2 \mathbf{U} - \mathbf{r}_i\mathbf{r}_i)\)，即可推导出公式 (6)。
\end{proof}

\section{动能 (Kinetic Energy)}

\begin{mydefinition}{惯量系数}
设 \(S\) 的构型由 \(n\) 个广义坐标表征，并定义了 \(n\) 个广义速度。那么 \(S\) 在 \(A\) 中的动能 \(K\) 可以表示为：
\begin{equation}
K = K_0 + K_1 + K_2 \tag{1}
\end{equation}
其中 \(K_j\) 是关于变量 \(u_1,\dots,u_n\) 的 \(j\) 次齐次多项式。特别地：
\begin{equation}
K_2 = \frac{1}{2}\sum_{r=1}^n \sum_{s=1}^n M_{rs} u_r u_s \tag{2}
\end{equation}
\(M_{rs}\) 称为**惯量系数**，定义为：
\begin{equation}
M_{rs} = \sum_{i=1}^N m_i \mathbf{v}_r^{P_i} \cdot \mathbf{v}_s^{P_i} \qquad (r,s=1,\dots,n) \tag{3}
\end{equation}
\end{mydefinition}

如果 \(B\) 是 \(S\) 中的刚体，其对 \(M_{rs}\) 的贡献为：
\begin{equation}
(M_{rs})_B = m\mathbf{v}_r \cdot \mathbf{v}_s + \boldsymbol{\omega}_r \cdot \mathbf{I} \cdot \boldsymbol{\omega}_s \qquad (r,s=1,\dots,n) \tag{4}
\end{equation}
其中 \(m\) 是 \(B\) 的质量，\(\mathbf{v}_r\) 是 \(B\) 质心的第 \(r\) 个偏速度，\(\boldsymbol{\omega}_r\) 是 \(B\) 的第 \(r\) 个偏角速度。

当广义速度选为广义坐标对时间的导数（即 \(u_r = \dot{q}_r\)）时，广义惯性力可以通过动能 \(K\) 表示为：
\begin{equation}
F_r^* = \frac{d}{dt}\left(\frac{\partial K}{\partial \dot{q}_r}\right) - \frac{\partial K}{\partial q_r} \tag{6}
\end{equation}

\begin{proof}[惯量系数与动能求导的推导]
将 \(K_2\) 的展开式与 \(M_{rs}\) 的定义进行比较即可得到 (3) 和 (4)。式 (6) 的证明利用了加速度与速度偏导数之间的关系：\(\mathbf{v}_r = \frac{\partial \mathbf{v}}{\partial \dot{q}_r}\) 以及 \(\mathbf{a} \cdot \mathbf{v}_r = \frac{d}{dt}(\mathbf{v} \cdot \frac{\partial \mathbf{v}}{\partial \dot{q}_r}) - \mathbf{v} \cdot \frac{\partial \mathbf{v}}{\partial q_r}\)，代入惯性力 \(F_r^* = \sum -m_i \mathbf{a}_i \cdot \mathbf{v}_r^{P_i}\) 后即可消去速度展开项，最终整理得到 (6)。
\end{proof}

\section{动力学方程 (Dynamical Equations)}

\begin{myTheorem}{Kane 动力学方程}
给定一个在牛顿参考系 \(N\) 中拥有 \(n\) 个自由度的系统 \(S\)，令 \(u_1,\dots,u_n\) 为 \(S\) 在 \(N\) 中的广义速度，\(F_1,\dots,F_n\) 为对应的广义主动力（4.1节），\(F_1^*, \dots, F_n^*\) 为对应的广义惯性力（4.3节）。
那么，\(S\) 的所有运动都满足以下方程：
\begin{equation}
F_r + F_r^* = 0 \qquad (r=1,\dots,n) \tag{1}
\end{equation}
这些方程称为 Kane 动力学方程。
\end{myTheorem}

\begin{proof}[Kane 方程的推导]
将系统视为由 \(\nu\) 个粒子组成。对每个粒子 \(P_i\)，根据牛顿第二定律有 \(\mathbf{R}_i = m_i \mathbf{a}_i\)。定义惯性力 \(\mathbf{R}_i^* = -m_i \mathbf{a}_i\)，则 \(\mathbf{R}_i + \mathbf{R}_i^* = 0\)。将该矢量等式与第 \(r\) 个偏速度 \(\mathbf{v}_r^{P_i}\) 点积，并对所有粒子求和：
\begin{equation}
\sum_{i=1}^\nu \mathbf{v}_r^{P_i} \cdot \mathbf{R}_i + \sum_{i=1}^\nu \mathbf{v}_r^{P_i} \cdot \mathbf{R}_i^* = 0
\end{equation}
根据广义主动力和广义惯性力的定义，上式等价于 \(F_r + F_r^* = 0\)。
\end{proof}

\section{线性化动力学方程 (Linearized Dynamical Equations)}

在为分析特定运动的稳定性或设计控制系统而推导动力学方程时，通常只需要线性化方程，即忽略 \(q_i\) 和 \(u_i\) 的二阶及更高阶项。这类方程可以不先写出精确方程，而按以下步骤进行：
1. 写出角速度、速度的完全非线性表达式，并求出偏角速度和偏速度；
2. 将所有角速度、速度和偏速度（及偏角速度）线性化；
3. 利用这些线性化后的量构建广义主动力和广义惯性力，并忽略所有二阶及以上项；
4. 代入 Kane 动力学方程 \(F_r + F_r^* = 0\) 得到所需方程。

\section{计算机化符号操作 (Computerization of Symbol Manipulation)}

当系统包含不止几个刚体或质点时，手动推导动力学方程极为繁重。此时，利用计算机进行文字形式的微分、乘法和代换显得尤为有利。当最终目标为微分方程的数值解时，可以通过计算机化符号操作直接生成数值求解的仿真程序代码，从而绕开手动书写显式的解析方程。

\section{离散多自由度系统 (Discrete Multi-Degree-of-Freedom Systems)}

当系统自由度很大时，数值求解微分方程可能非常耗时。通过引入相关经典振动问题的**模态坐标**，并排除部分高频模态，可以显著减少计算量。

\begin{myTheorem}{模态分析与降阶}
对于描述微小振动的常数矩阵线性微分方程：
\begin{equation}
M\ddot{\mathbf{x}} + S\mathbf{x} = 0 \tag{1}
\end{equation}
其中 \(M\) 和 \(S\) 是 \(n \times n\) 矩阵。可通过引入一个 \(n \times \nu\)（\(\nu < n\)）的常数模态矩阵 \(A\) 和一个 \(\nu \times 1\) 的时间相关矩阵 \(\mathbf{q}(t)\) 来近似 \(\mathbf{x}\)：
\begin{equation}
\mathbf{x} = A\mathbf{q} \tag{2}
\end{equation}
\end{myTheorem}

构造模态矩阵 \(A\) 的步骤如下：
1. 求出 \(M^{-1}S\) 的前 \(\nu\) 个非零特征值 \(\lambda_1,\dots,\lambda_\nu\)；
2. 计算对应的特征向量 \(\mathbf{B}_1,\dots,\mathbf{B}_\nu\)；
3. 根据质量矩阵 \(M\) 进行归一化，得到模态矩阵 \(A = [\mathbf{A}_1\ \mathbf{A}_2\ \dots\ \mathbf{A}_\nu]\)。

矩阵 \(A\) 满足正交性条件：
\begin{equation}
A^T M A = U, \quad A^T S A = \Lambda \tag{11}
\end{equation}
其中 \(U\) 为 \(\nu \times \nu\) 单位矩阵，\(\Lambda\) 是以特征值 \(\lambda_i\) 为对角元的对角矩阵。
系统的初始条件可通过伪逆求得：
\begin{equation}
\mathbf{q}(0) = A^T M \mathbf{x}(0), \quad \dot{\mathbf{q}}(0) = A^T M \dot{\mathbf{x}}(0) \tag{6}
\end{equation}

\section{航天器的集中质量模型 (Lumped Mass Models of Spacecraft)}

某些航天器可以被建模为通过弹性连接的一组质点。例如，无约束的桁架结构的运动可以通过将质点放置在桁架节点上来分析，每个质点的质量等于连接在该节点上的所有杆件质量的一半，用无质量的弹簧来表示杆件的刚度属性。
当结构单元数量庞大时，数值积分过程耗时巨大，此时可再次采用 4.8 节的模态排除技术。通过构建结构的质量矩阵 \(M\) 和刚度矩阵 \(S\)，借助计算机程序求解特征值问题 \(M^{-1}S\) 的模态矩阵 \(A\)，将 \(3n\) 个节点的位移 \(\mathbf{x}\) 分解为刚体平动、转动与少量弹性模态的叠加，进而高效求解。
\newpage
\section{具有连续弹性组件的航天器 (Spacecraft with Continuous Elastic Components)}

当航天器包含由刚体支撑的连续弹性组件时，相对于刚体的小位移运动通常由无法用分离变量法求解的偏微分方程控制。但若支撑刚体固定于牛顿参考系中（如悬臂梁），则可用分离变量法求解，进而可基于这些振动模态建立近似动力学方程。

\begin{mydefinition}{欧拉-伯努利梁方程}
对于长度为 \(L\)、抗弯刚度为 \(EI\)、单位长度质量为 \(\rho\) 的均匀悬臂梁，其横向小挠度振动由下式控制：
\begin{equation}
EI\frac{\partial^4 y}{\partial x^4} + \rho\frac{\partial^2 y}{\partial t^2} = 0 \tag{1}
\end{equation}
满足边界条件 \(y(0,t)=y'(0,t)=y''(L,t)=y'''(L,t)=0\)。
\end{mydefinition}

其通解可表示为模态叠加形式：
\begin{equation}
y = \sum_{i=1}^{\infty} \phi_i q_i \tag{3}
\end{equation}
其中模态函数 \(\phi_i\) 和模态坐标 \(q_i\) 分别为：
\begin{equation}
\phi_i = \cosh\frac{\lambda_i x}{L} - \cos\frac{\lambda_i x}{L} - \frac{\cosh\lambda_i + \cos\lambda_i}{\sinh\lambda_i + \sin\lambda_i}\left( \sinh\frac{\lambda_i x}{L} - \sin\frac{\lambda_i x}{L} \right) \tag{4}
\end{equation}
\begin{equation}
q_i = a_i \cos p_i t + \beta_i \sin p_i t \tag{5}
\end{equation}
其中 \(\lambda_i\) 是超越方程 \(\cos\lambda \cosh\lambda + 1 = 0\) 的根，固有频率为 \(p_i = \lambda_i^2 \sqrt{\frac{EI}{\rho L^4}}\)。
模态函数满足正交性条件：
\begin{equation}
\int_0^L \phi_i \phi_j \rho dx = m_B \delta_{ij} \tag{8}
\end{equation}
\begin{equation}
\int_0^L EI \phi_i'' \phi_j'' dx = p_i^2 m_B \delta_{ij} \tag{9}
\end{equation}

\section{用有限元法构建模态函数 (Use of the Finite Element Method for the Construction of Modal Functions)}

当组件的抗弯刚度 \(EI\) 或线密度 \(\rho\) 不是常数，导致无法解析求解偏微分方程时，可采用有限元法构造近似的模态函数。

\begin{mydefinition}{梁单元的质量与刚度矩阵}
将梁划分为 \(m\) 个棱柱单元，第 \(i\) 个单元长度为 \(L_i\)，抗弯刚度为 \((EI)_i\)，线密度为 \(\rho_i\)。每个单元利用端部位移 \(z_{2i-1}, z_{2i}\) 和转角 \(z_{2i+1}, z_{2i+2}\) 作为节点自由度，定义形状函数 \(\psi_{i1},\dots,\psi_{i4}\)。
该梁单元的质量矩阵 \(M_i\) 和刚度矩阵 \(S_i\) 分别为：
\begin{equation}
M_i = \frac{\rho_i L_i}{420}\begin{bmatrix} 156 & 22L_i & 54 & -13L_i \\ 22L_i & 4L_i^2 & 13L_i & -3L_i^2 \\ 54 & 13L_i & 156 & -22L_i \\ -13L_i & -3L_i^2 & -22L_i & 4L_i^2 \end{bmatrix} \tag{12}
\end{equation}
\begin{equation}
S_i = \frac{(EI)_i}{L_i^3}\begin{bmatrix} 12 & 6L_i & -12 & 6L_i \\ 6L_i & 4L_i^2 & -6L_i & 2L_i^2 \\ -12 & -6L_i & 12 & -6L_i \\ 6L_i & 2L_i^2 & -6L_i & 4L_i^2 \end{bmatrix} \tag{13}
\end{equation}
\end{mydefinition}

将各单元的质量矩阵 \(M_i\) 和刚度矩阵 \(S_i\) 组装成全局质量矩阵 \(M\) 和全局刚度矩阵 \(S\)，即可对无约束梁（自由-自由）的振动进行模态分析。

当梁受到约束时（如一端固定），则从全局矩阵中删除对应的自由度行和列，形成受约束梁的质量矩阵 \(\overline{M}\) 和刚度矩阵 \(\overline{S}\)，并求解其对应的特征值 \(\lambda_i\) 和特征向量。

若将有限元求得的近似模态函数 \(\Phi_i\) 代入 4.10 节的方程中，就可以对具有非均匀特性的连续弹性组件进行建模并构建整个航天器的动力学方程。